% =====================================================================
% Track 4 (Sprint 1) — §III.23 Symmetry projection / Young tableau
% Draft for inclusion in Paper 34 §III, after §III.22.
%
% Three-axis dictionary row (for §IV Table~\ref{tab:projection_axes}):
%
%   Projection            : Symmetry / Young tableau
%   Variables added       : irrep tableau $\lambda$
%                           (e.g.\ $[1^N]$ antisymmetric, $[2,2]$ singlet,
%                           $[3,1]$ mixed)
%   Dimension             : preserved (orthogonal projection on the
%                           $N$-particle Hilbert space)
%   Transcendental signature : \textit{none} — character coefficients
%                           are integers, projector matrix entries are
%                           rationals
%   Layer transition      : Layer 1 $\to$ Layer 1 (purely combinatorial
%                           projection on the $N$-particle Hilbert
%                           space; no transcendental injection)
%
% Cross-references: \S~\ref{sec:proj_spinor} (Camporesi--Higuchi spinor
% bundle: the single-particle spin structure that the multi-particle
% $S_N$ projection composes with), Paper~16 ($S_N$ representation
% theory framework underlying the atomic classifier), Paper~17 §VIII
% (Level~4N $S_4$ $[2,2]$ projection for 4-electron LiH), Paper~27
% (entanglement structure of antisymmetrized states).
% =====================================================================

\subsection{Symmetry projection / Young tableau}
\label{sec:proj_symmetry_tableau}
\textbf{Source:} the unrestricted $N$-particle tensor product
$\mathcal{H}^{\otimes N}$ on the single-particle Fock-projected $S^3$
graph (or on any single-particle natural-geometry sector;
\S~\ref{sec:proj_fock}, \S~\ref{sec:proj_bargmann},
\S~\ref{sec:proj_molframe}).  At this stage indistinguishability of
identical fermions has not yet been enforced; the Hilbert space
contains symmetric, antisymmetric, and mixed-symmetry components
without distinction.\\
\textbf{Target:} the $\lambda$-irrep subspace
$\mathcal{H}_\lambda \subset \mathcal{H}^{\otimes N}$, obtained via
the character-based orthogonal projector
\[
  P_\lambda \;=\; \frac{d_\lambda}{|S_N|}\,\sum_{g \in S_N}
                  \chi_\lambda(g)\,\rho(g),
\]
where $\lambda$ is a Young tableau (partition of $N$), $d_\lambda$ is
the dimension of the corresponding irreducible representation of
$S_N$, $\chi_\lambda(g)$ is its character at the permutation $g$, and
$\rho(g)$ is the natural action of $g$ on $\mathcal{H}^{\otimes N}$
by index permutation.  For fully indistinguishable spin-$\tfrac{1}{2}$
fermions in the unrestricted formulation, $\lambda = [1^N]$
(antisymmetric); in the Slater-determinant / spatial-spin coupled
formulation, $\lambda$ is the conjugate of the spin Young diagram
(Paper~16 §\ref{sec:sn_irreps} convention), so that $S_N$
antisymmetry is implemented \emph{jointly} on spatial and spin
indices.\\
\textbf{Variables introduced:} the irrep tableau $\lambda$ itself
(a discrete combinatorial label; equivalently a partition of $N$).
For $N=2$: $\lambda \in \{[2],\,[1,1]\}$ (triplet vs singlet under
spatial-spin coupling).  For $N=4$ closed shell: $\lambda = [2,2]$
is the spin-singlet sector (Paper~17 §VIII, Track AJ).  For general
$N$ with total spin $S$: $\lambda_1 = N/2 + S$ and
$\lambda_2 = N/2 - S$ (Paper~16 §\ref{sec:sn_irreps}).\\
\textbf{Dimension introduced:} \emph{none}.  $P_\lambda$ is an
orthogonal projection on the $N$-particle Hilbert space; it does not
introduce a new physical scale, focal length, or coordinate.  The
projector reduces the linear dimension of the working space from
$(\dim\mathcal{H})^N$ to a $\lambda$-irrep-counted subspace, but no
new $[L]$, $[E]$, or $[T]$ enters the framework.\\
\textbf{Transcendental signature:} \textit{none}.  The character
table of $S_N$ is integer-valued (Frobenius character formula); the
projector $P_\lambda$ has rational matrix entries (denominators
dividing $|S_N| = N!$, numerators integer); the dimension $d_\lambda$
is integer (hook-length formula).  No $\pi$, no $\zeta$, no Hurwitz
content enters at this step.  The projection is the unique
\S~\ref{sec:layer2} entry whose entire algebraic content lives in
$\mathbb{Z}/N!$, making it the cleanest example of a
\emph{combinatorial} Layer~1 $\to$ Layer~1 projection.\\
\textbf{Used in:} all $N \geq 2$ computations in the framework, at
multiple levels of explicitness:
\begin{itemize}
\item \textbf{Level 3 (He, $N=2$).}  Singlet $[1,1]$ vs triplet $[2]$
spatial-spin coupling is the projection that splits the He $1^1 S$
ground state from the $2^3 S$ first excited state.  Spatial-spin
coupling for two electrons with total spin $S \in \{0, 1\}$ yields
$\lambda_\text{spatial} = [1,1]$ for $S=0$ (antisymmetric spatial,
symmetric spin: singlet) and $\lambda_\text{spatial} = [2]$ for
$S=1$ (symmetric spatial, antisymmetric spin: triplet).
\item \textbf{Level 4N (LiH, $N=4$).}  The $S_4$ $[2,2]$ Young
diagram projection is the character-based projector onto the
closed-shell singlet sector of the 4-electron mol-frame
hyperspherical Hilbert space on $S^{11}$ (Paper~17 §VIII; Track AJ,
\texttt{geovac/n\_electron\_solver.py}).  The projection reduces the
raw SO(12) angular channel count by a tableau-dependent factor and
makes the 4-electron PES computationally tractable; it is empty at
$l_{\max} = 0$ (no antisymmetric 4-electron state can be built from
$s$-waves), gives 2 channels at $l_{\max} = 1$, 12 channels at
$l_{\max} = 2$, and 40 channels at $l_{\max} = 3$.
\item \textbf{Graph-native CI (all $N$).}  The Slater-determinant
basis used in \texttt{geovac/casimir\_ci.py} is itself the
\emph{Slater-determinant realization} of the $[1^N]$ antisymmetric
projector: each determinant
$\det[\phi_{i_1}(\vec{r}_1)\cdots\phi_{i_N}(\vec{r}_N)]$ is the
image under $P_{[1^N]}$ of the single Hartree product
$\phi_{i_1}(\vec{r}_1)\otimes\cdots\otimes\phi_{i_N}(\vec{r}_N)$,
up to a normalization factor.
\item \textbf{Composed atomic classifier (Paper~16).}  The atomic
structure-type assignment (A/B/C/D/E) for $Z = 1$--$30$ is itself
read off the $S_N$ irrep shape: structure-type letters correspond to
specific tableau patterns under the periodic-table indexing of
Paper~16 §\ref{sec:sn_irreps}.  The classifier is the boundary at
which the $S_N$ projection becomes computationally explicit in the
production composed-Hamiltonian pipeline.
\item \textbf{Hyperfine and cross-register couplings (Track NI,
Paper~23 §VI).}  In the composed nuclear--electronic deuterium
Hamiltonian, the Clebsch--Gordan multiplicity $(2I+1)/2$ relating
$F = I+J$ to $F = I-J$ hyperfine levels (D HFS autopsy,
\S\ref{sec:autopsy_d_hfs}) is itself a small-$N$ tableau projection
applied to a coupled spin Hilbert space.  See \S~\ref{sec:proj_3j}
for the angular-momentum coupling content; the present projection
is its $S_N$ analog on the indistinguishable-particle side.
\end{itemize}
\textbf{Empirical anchor:} He $2^1 S$--$2^3 S$ exchange splitting at
$6421.46$~cm$^{-1}$ (NIST experimental) is the cleanest finite-energy
witness of the singlet-vs-triplet ($[1,1]$ vs $[2]$ spatial)
projection in the catalogue.  The splitting is not a small Layer-2
correction; it is the energy difference between two distinct
$\lambda$-irrep sectors of the same single-particle basis, generated
entirely by the exchange-integral content of the two-electron $V_{ee}$
matrix elements under spatial antisymmetrization vs symmetrization.
Sprint Calc-He-2S-2S-splitting and Sprint Hylleraas-$r_{12}$
(2026-05-09) returned this entry at $+11.86\%$ (graph-native CI,
$n_{\max} \leq 11$) and $+209\%$ (single-$\alpha$ Hylleraas-3p,
$\omega = 5$) respectively, with the diagnosis that the residuals
trace to single-exponent basis limitations for the excited
$\lambda$-irrep state rather than to the projection itself
(\S~\ref{sec:offprecision} He singlet-triplet row).  At the
projection level, the splitting is structurally clean: the projector
$P_{[1,1]}$ vs $P_{[2]}$ acts identically on either side; only the
$V_{ee}$ matrix elements respond differently to the two sectors.\\
\textbf{Honest scope.}  At modest $N$ ($N \leq 4$) the framework
implements the character-based projector $P_\lambda$ directly:
explicit construction of the projector matrix in the Slater /
configuration-state-function basis, or via the equivalent Young
symmetrizer in the unrestricted tensor product.  This is exact and
combinatorial.  At large $N$ (heavy atoms, $N \gg 10$), the framework
falls back on \emph{compositional approximations} to the tableau
projection: (i) single-Slater-determinant Hartree--Fock reductions
project onto a single representative of the $[1^N]$ irrep rather
than spanning it, (ii) frozen-core constructions
(\texttt{geovac/neon\_core.py} for [Ne]-core, [Kr]-core, [Xe]-core
species) project the antisymmetric structure onto a product of a
fixed core determinant with a smaller-$N$ valence tableau projection,
(iii) the composed-block architecture for multi-center systems
(\S~\ref{sec:proj_molframe}) treats each block's antisymmetric
subspace independently, neglecting inter-block antisymmetry except
through cross-block ERI / cross-center $V_\text{ne}$ corrections.
Each of these is itself a structured projection onto a coarser
tableau pattern.  The boundary between exact character-projection
and compositional approximation is empirical, set by the largest $N$
at which $|S_N| = N!$ direct construction remains tractable.  As of
v2.40.0 this boundary sits at $N \approx 6$ for the atomic
classifier's explicit construction (extended via lookup tables to
higher $N$ where the periodic-table structure-type pattern repeats),
and at $N = 4$ for the Level~4N solver
(\texttt{geovac/n\_electron\_solver.py}).\\
\textbf{Structural reading.}  This is the projection at which
\textbf{Pauli antisymmetry enters the framework explicitly} rather
than implicitly.  In many quantum-chemistry frameworks
antisymmetrization is folded into the coordinate or operator choice
(e.g.\ second-quantized anticommutators built into the Hamiltonian's
algebraic structure); in GeoVac the dimensionless graph
(\S~\ref{sec:bare_graph}) is symmetry-agnostic by construction, and
Pauli antisymmetry is recovered as a separate, named, character-based
projection at the multi-particle level.  This is the structural
sibling of \S~\ref{sec:proj_spinor} (which introduces the
spin-$\tfrac{1}{2}$ Hilbert space at the single-particle level): the
spinor projection commits to particle statistics (fermionic
half-integer spin), while the present projection commits to the
many-body consequence of those statistics ($N$-particle
antisymmetry).  Composition $P_\lambda \circ
\text{spinor\_lift}$ realizes the standard spatial-spin coupling
scheme; composition $P_\lambda \circ \text{Fock}$ realizes the
$N$-fermion graph Hilbert space on which graph-native CI operates.
Categorically distinct from \S~\ref{sec:proj_3j} (angular-momentum
coupling, $SU(2)$ Clebsch--Gordan, continuous-symmetry character
content) and \S~\ref{sec:proj_wignerD} (rotation between rotated
coordinate frames, $SO(3)$ representation theory).  The shared
abstraction across these three is that all three are character-based
projections onto irreducible representations of finite or compact
groups acting on the Hilbert space; the structural distinction is
the group: $S_N$ for the present projection, $SU(2)$ for
\S~\ref{sec:proj_3j}, $SO(3)$ for \S~\ref{sec:proj_wignerD}.  The
$S_N$ character table has integer entries while the $SU(2)$ and
$SO(3)$ projector matrix elements introduce algebraic-extension and
trigonometric content respectively (\S~\ref{sec:dictionary}
Table~\ref{tab:projection_axes}); the absence of transcendental
content in the symmetry projection is therefore not coincidental
but reflects the discrete finite-group structure underlying Pauli
exclusion.
